\documentclass[10pt,a4paper]{amsart}
\usepackage[margin=19mm]{geometry}
\usepackage{amsmath,amssymb,mathtools}
\usepackage[scr=rsfs,cal=euler]{mathalfa}
\usepackage{xfrac}
\usepackage[unicode,hidelinks,bookmarks=false]{hyperref}
\newcommand{\ie}{i.e.\ }
\newcommand{\cf}{cf.\ }
\newcommand{\resp}{respectively\ }
\newcommand{\Subsection}{\subsection}
\providecommand{\nobreakdash}{\nobreak\hbox{-}}
\setlength{\emergencystretch}{2em}
\multlinegap0pt
\usepackage{amssymb}
\usepackage[all]{xy}
\usepackage[shortlabels]{enumitem}
\usepackage{bbm}
\usepackage{bm}
\usepackage{mathtools}
\usepackage{multirow}
\newtheorem{corollary}{Corollary}
\newcommand{\bP}{\mathbb{P}}
\newcommand{\bQ}{\mathbb{Q}}
\newcommand{\calO}{\mathcal{O}}
\title{A guide to wall crossing for moduli of varieties}
\author{Kristin DeVleming}
\date{Source: arXiv:2602.20286v1 (CC BY 4.0)}
\begin{document}
\maketitle

\noindent\emph{Source and licence.} A guide to wall crossing for moduli of varieties. Version arXiv:2602.20286v1 (CC BY 4.0), distributed by arXiv under the Creative Commons Attribution 4.0 International licence, http://creativecommons.org/licenses/by/4.0/, as recorded in the arXiv licence metadata for this exact version and shown on the abstract page https://arxiv.org/abs/2602.20286v1. Full licence notice: \texttt{licenses/arxiv-2602.20286-CC-BY-4.0.txt}.\par\smallskip

\noindent\textit{Editorial scope.} Counted unit: the article's corollary corollary (source0.tex lines 1574--1576). The article's complete original proof of it is delivered with it (source0.tex lines 1578--1596, one \texttt{proof} environment of 19 lines). No mathematics is written, repaired or re-derived: the statement and the proof are the article's own lines, in the article's own notation, and no retained line is substituted or re-flowed.  Retained uncounted as the source-local context the proof needs: lines 1514--1518.  Nothing was omitted from any retained span, so the package carries no omission record.  No retained line contains a citation, so the wrapper carries no bibliography and no bibliography entry.  No retained line contains a figure environment, an image-inclusion command or drawing code, so no image is delivered and the package declares no asset dependency.  Editorial formatting only, in addition: an AMS \texttt{amsart} preamble that restates the article's own declarations and macros used by the retained lines, namely the environments \texttt{corollary} on separate counters, together with the standard packages those lines need.  The retained environments print in this document as Corollary 1 in order of appearance, whereas the article prints the same items with the numbering of its own full document.  The complete native source of the article as distributed in the arXiv e-print for this exact version is bundled unchanged as \texttt{sources/195174/source0.tex}.
\begin{corollary}\label{cor:manycomps}
    Let $X$ be $\bQ$-Gorenstein slc surface with $k$ components and $D \subset X$ a curve with $j$ components and an $A_n$ singularity satisfying the conditions in the previous theorem.  Then, for any $c \in (\frac{1}{2} + \frac{1}{n+1}, 1]$ and $\ell$ such that $c <\frac{1}{2} + \frac{1}{\ell+1}$, there exists a specialization $Y$ of $X$ with $k+n -\ell$ components containing a divisor $D_Y$ with $j + n - \ell$ components such that $D_Y$ has an $A_\ell$ singularity and $(Y, c D_Y)$ is a stable pair.

    If $c = 1$, and $\ell = 1$, this produces a stable pair $(Y,D_Y)$ with $k+n -1$ components. 
\end{corollary}

\begin{corollary}
    In the KSBA moduli space compactifying the locus pairs $(\bP^2, D)$ where $D \in |\calO(d)|$ and $d \ge 6$, there exists a surface with at least $d^2/2 - 1$ components if $d$ is even and $d(d-1)/2  - 1$ components if $d$ is odd.  
\end{corollary}

\begin{proof}
    First, assume $d$ is even.  We construct a curve of degree $d$ with an $A_n$ singularity where $n = d^2/2 - 1$.  To do this, consider the plane curve $C$ given by $f = 0$, where
    \[ f(x,y,z) = (xz^{d/2-1} - y^{d/2})^2 - x^d. \]
    By direct computation, this curve has an isolated singular point at $[0:0: 1]$.  In the coordinate chart where $z \ne 0$, we may write 
    \[ f(x,y,1) = (x - y^{d/2})^2 - x^d\]
    and consider the change of coordinates given by $x' = x - y^{d/2}$, so this becomes 
    \[ (x')^2 - (x'+y^{d/2})^d = (x')^2 - y^{d^2/2} + h.o.t.\]
    where the higher order terms are with respect to the weighting of $(x', y)$ by $(d^2/2, 2)$.  This defines an $A_n$ singularity where $n = d^2/2 - 1$.  By construction, this curve is the union of two smooth degree $d/2$ curves as $f$ factors as 
    \[ f(x,y,z) = (xz^{d/2-1} - y^{d/2})^2 - x^d = (xz^{d/2-1} - y^{d/2} +x^{d/2})(xz^{d/2-1} - y^{d/2} -x^{d/2}).\]
    Because $d/2 \ge 6/2 = 3$, the genus of each component of $C$ is at least one. 

    When $d$ is odd, we construct a plane curve with an $A_n$ singularity where $n = d(d-1)/2  - 1$ in the same way.  This curve $C$ is given by the vanishing of the equation 
    \[ f(x,y,z) = (xz^{(d-1)/2-1} - y^{(d-1)/2})^2z - x^d. \]
    As above, this has an $A_n$ singularity at $[0:0:1]$ as claimed.  This curve is irreducible and its geometric genus is  
    \[ \frac{(d-1)(d-2)}{2} - \lfloor \frac{d(d-1)}{4}\rfloor \ge \frac{(d-1)(d-2)}{2} - \frac{d(d-1)}{4} = \frac{(d-1)(d-4)}{4} > 0  \]
    as $d \ge 6$.

    For any $c > \frac{3}{d} $, $K_{\bP^2} + cC$ is ample, and for any degree $d \ge 6$, $c = \frac{1}{2} + \frac{1}{n} > \frac{1}{2} \ge \frac{3}{d}$, so the hypothesis $K_{\bP^2} + cC$ is ample in the previous theorem is satisfied.  The hypotheses of Corollary \ref{cor:manycomps} are therefore all satisfied, so the statement holds. 
\end{proof}

\end{document}
